\begin{table}[t]
\centering
\caption{Signal ablation across the four independent validation runs.}
\label{tab:ablation}
\begin{tabular}{lllllcp{3.2cm}}
\toprule
Detector & Gaussian attribution & Gaussian onset err/delay & Directional attribution & Directional onset err/delay & Honest persistent FP (4 runs) & Interpretation \\
\midrule
Peer-only & No (0/2) & wrong client & No (0/2) & wrong client & 0 & Never contributes a detection in any development or validation run \\
Magnitude-only & Yes (2/2) & err 0/delay 2 & Yes (2/2) & seed123: err 27/delay 29; seed2026: err 24/delay 26 & 0 & Fast/correct for Gaussian; correct but far too late for directional \\
Temporal-residual-only & Yes (2/2) & seed123: err 1/delay 3; seed2026: err 0/delay 2 & Yes (2/2) & err 0/delay 2 & 0 & Fast/correct for directional; matches or nearly matches full detector for Gaussian \\
Full FedTrace v3 (magnitude + temporal-residual) & Yes (2/2) & err 0/delay 2 & Yes (2/2) & err 0/delay 2 & 0 & Matches whichever single signal is informative -- coverage, not fusion \\
\bottomrule
\end{tabular}
\vspace{2pt}\par\raggedright\footnotesize Message: the two signals provide complementary attack-family coverage. This table does NOT support, and does not claim, that the combiner improves over the best single detector -- in every run checked, the full-detector outcome matches whichever single component was already informative for that run.
\end{table}
